This appendix is provided in compliance with Section~IV.b of the De~La~Salle University \textit{Policy for Generative AI in Education}. It discloses the generative AI tools used during the preparation of this thesis, the extent to which their outputs contributed to the overall manuscript and system, and the verification steps performed by the researchers on any material derived, in part or in whole, from these tools.

\section*{Generative AI Use Disclosure}

\textbf{ChatGPT} (OpenAI) was used to improve the grammar, clarity, and professional tone of selected portions of the manuscript. It was also used to assist in planning the thesis timeline and organizing tasks throughout the study to support efficient project management, and -- alongside Claude -- to assist with recurring \LaTeX{} formatting tasks (tables, figures, lists, and other repeated formatting fixes) so that the manuscript would render with consistent visual formatting. It was not used to generate citations or to introduce any source into the reference list. The researchers used the free tier of ChatGPT; no paid subscription was utilized.

\textbf{NotebookLM} (Google) was used only to organize and summarize literature that the researchers had already located, retrieved, and uploaded themselves. NotebookLM cannot generate outbound links to publisher pages, DOIs, or repositories: it can only cite the sources that the researchers themselves have supplied to it, and it will not answer any question about a source until the researchers have first placed that source into the notebook. Every hyperlink, DOI, and bibliographic identifier in this manuscript was therefore researched, obtained, and typed in by the researchers themselves from the primary source (IEEE~Xplore, ACM~Digital~Library, ScienceDirect, MDPI, Springer, arXiv, publisher web pages, and the corresponding DOI resolvers at \url{https://doi.org/}). The full sourcing and verification workflow that the researchers followed before any reference was admitted into the bibliography is as follows:
\begin{enumerate}
    \item Performed the same Review of Related Literature (RRL) sourcing procedure that had been established during REMETHS: repeated keyword searches across multiple academic search engines -- primarily Google~Scholar, IEEE~Xplore, ScienceDirect, ACM~Digital~Library, and Springer -- using the tag-line queries developed for this study (fuel monitoring, GPS tracking, IoT fleet telematics, fuel-consumption anomaly detection, isolation forest, matrix profile, and related terms) in order to surface studies with subject matter similar to this thesis.
    \item Re-executed the same sourcing procedure a second time during PRO1, this time following the PRISMA identification--screening--eligibility--inclusion flow for the methodology chapter, so that the pool of studies drawn on for the methodology was reproducible and independent of the initial RRL pass rather than the ad-hoc RRL pass alone.
    \item Assessed the relevance of each candidate study to the objectives of this thesis, discarding studies that did not directly inform the fuel-monitoring, GPS-tracking, or anomaly-detection scope.
    \item Cross-verified each retained study by inspecting its own reference list and citation trail on the publisher-of-record record, to confirm that the work was cited as attributed and had not been misidentified.
    \item Verified the title, venue, year, volume/issue, and page range against the publisher-of-record record, and recorded the DOI (or, where no DOI exists, the canonical publisher/repository URL) directly in the BibTeX entry so that every claim in the manuscript resolves to an accessible primary source.
    \item Only after the source had been retrieved and validated by the above steps was it uploaded into NotebookLM, which was then used to extract specific passages and produce section-level summaries for easier note-taking; because NotebookLM cites only the sources that the researchers themselves have provided, the tool at no point introduced a citation that the researchers had not already retrieved and verified.
\end{enumerate}
In short, the researchers themselves \textbf{accessed, verified, and cited} every source: NotebookLM was used only as an organizational and note-taking aid over material the researchers had already retrieved and verified through the RRL and PRISMA sourcing procedures above, and it neither discovered nor introduced any of the citations in this manuscript. The researchers used the free tier of NotebookLM; no paid subscription was utilized.

\textbf{Claude} (Anthropic) was used for technical brainstorming, code outlining, and debugging assistance when the researchers encountered implementation difficulties. In addition, Claude was used -- alongside ChatGPT -- to assist with \LaTeX{} formatting tasks (table layouts, figure environments, list structures, bibliography-macro adjustments, and repeated formatting fixes) so that the manuscript would compile cleanly and render with consistent visual formatting. All suggested code, \LaTeX{} snippets, and solutions were manually reviewed, modified when necessary, tested against the actual embedded, backend, machine-learning, and \LaTeX{} build environments, and validated by the researchers before being incorporated into the system or the manuscript. Claude was not used to generate citations, DOIs, or bibliographic entries. The researchers used a paid Claude~Pro subscription to obtain extended session limits during the implementation and debugging phases of the project.
